\section{Introducción}

El presente documento desarrolla el tema \textit{Análisis financiero: análisis horizontal y análisis vertical de estados financieros}, correspondiente al curso Contabilidad Empresarial de la Universidad Nacional del Callao. El contenido se alinea con la Unidad IV: análisis e interpretación de estados financieros.

El análisis financiero permite transformar datos contables en información útil para evaluar la situación económica y financiera de una empresa. En ese proceso, los estados financieros no se leen como cifras aisladas, sino como información comparable, estructurada e interpretable.

La exposición se organiza en cuatro bloques: fundamentos, análisis horizontal, análisis vertical y aplicación empresarial. El caso Empresa ABC permite mostrar el uso de fórmulas, cuadros, gráficos e interpretación técnica.
